\chapter{Elektronika i łączność}\label{ch:elektronika}

\begin{sekcja}{Architektura}
Elektronika platformy jest w gondoli platformy. Zawiera:
\begin{itemize}
  \item \textbf{komputer bezpieczeństwa} (HABduino): GPS, telemetria, odcinacz główny,
  \item \textbf{komputer pokładowy} (OBC): zapis danych, zapasowe funkcje bezpieczeństwa, odcinacz zapasowy,
  \item \textbf{lokalizatory po lądowaniu}: SPOT Trace, dwa znaczniki Bluetooth (Apple i Google),
  \item baterie.
\end{itemize}

\begin{center}\resizebox{\linewidth}{!}{\rysTorRadiowy}\end{center}
\captionof{figure}{Tor radiowy: od gondoli przez stacje odbiorcze do map śledzenia.}\label{fig:tor-radiowy}
\end{sekcja}

\begin{sekcja}{Komputer bezpieczeństwa: HABduino}
HABduino to nakładka na Arduino do trackerów balonowych~\cite{habduino}. Nadaje telemetrię RTTY (moduł Radiometrix NTX2B, \SI{434}{MHz}) i APRS (\SI{144.8}{MHz}). Projekt nie jest już rozwijany. Zalecane oprogramowanie to FlexTrack~\cite{flextrack}: obsługuje HABduino~v4, RTTY i APRS, zawiera obsługę odcinacza. Opis budowy trackera: \cite{akerman}.

HABduino jest \textbf{głównym komputerem bezpieczeństwa}:
\begin{itemize}
  \item podaje pozycję (GPS) przez dwa niezależne łącza radiowe,
  \item steruje odcinaczem głównym (na komendę i automatycznie),
  \item działa niezależnie od OBC i od ładunku.
\end{itemize}
\end{sekcja}

\begin{sekcja}{Komputer pokładowy (OBC)}
OBC dodatkowo implementuje te same funkcje bezpieczeństwa, ale niezależnie od HABduino:
\begin{itemize}
  \item własny odbiornik GPS i własne kryteria zakończenia lotu,
  \item \textbf{niezależny odcinacz innego typu} (np. z drutem oporowym, gdy główny jest serwomechaniczny) -- jedna przyczyna awarii nie wyłącza obu,
  \item zapis danych z czujników na karcie SD, ze znacznikiem czasu,
  \item zapis stanu misji w pamięci nieulotnej: restart w locie nie zeruje misji.
\end{itemize}

\textbf{Wymagania:}
\begin{itemize}
  \item Wysokość bierz z GPS. Wzór na wysokość z ciśnienia jest niedokładny powyżej około \SI{10}{km}.
  \item Odrzucaj błędne odczyty czujników (kody błędów, wartości poza zakresem).
  \item Czujniki ciśnienia muszą mierzyć do około \SI{5}{hPa} (BME280 mierzy tylko do \SI{300}{hPa}).
  \item Dodaj diody statusu: zasilanie, GPS, karta SD, stan odcinacza.
\end{itemize}
Kod i dokumentacja OBC: \url{https://github.com/simle-stardust}.
\end{sekcja}

\begin{sekcja}{Konfiguracja trackera}
\begin{enumerate}
  \item Ustaw znak wywoławczy krótkofalowca (\emph{callsign}).
  \item Ustaw SSID APRS (zwyczajowo \texttt{-11} dla balonów).
  \item Ustaw identyfikator ładunku RTTY (\emph{payload ID}), taki sam jak w SondeHub.
  \item Ustaw moduł GPS w tryb \emph{airborne} (działa powyżej \SI{18}{km}; moduły u-blox do około \SI{50}{km}).
  \item Wgraj program. Zapisz wersję (numer commita albo wydania) w karcie misji.
\end{enumerate}

\begin{uwaga}Skonfiguruj tracker i serwisy śledzenia \textbf{najpóźniej 2 dni przed lotem}. Niektóre ustawienia (np. konfiguracja ładunku w SondeHub) potrzebują czasu, aby zadziałały. Dzień przed lotem zostaw na weryfikację, logistykę i obserwację pogody.\end{uwaga}
\end{sekcja}

\begin{sekcja}{Anteny i test zasięgu}
\begin{itemize}
  \item Antenę nadajnika montuj pod gondolą, skierowaną w dół.
  \item Antenę GPS montuj na górze gondoli, z widokiem na niebo, z dala od baterii i metalu.
  \item Po montażu sprawdź złącza i przewody anten.
\end{itemize}

\begin{uwaga}Z bliska odbiornik odbierze sygnał nawet bez anteny -- wystarczy promieniowanie z płytki i przewodów. Test obok stołu niczego nie dowodzi. Sprawdź odbiór z odległości \textbf{co najmniej \SI{100}{m}}, z anteną w konfiguracji lotnej (w zamkniętej gondoli).\end{uwaga}
\end{sekcja}

\begin{sekcja}{Lokalizatory po lądowaniu}
Przy ziemi sygnał z trackera zanika: zasłaniają go drzewa, budynki i ukształtowanie terenu. Lokalizatory pomagają znaleźć gondolę po lądowaniu.

\begin{itemize}
  \item \textbf{SPOT Trace} -- lokalizator satelitarny. Działa tylko do około \SI{6500}{m}, więc służy tylko po lądowaniu. Wymaga aktywnej subskrypcji. Ustaw interwał raportów (np. 5 minut). Antena do góry.
  \item \textbf{Znaczniki Bluetooth: Apple AirTag i znacznik sieci Google (Find Hub)} -- najlepiej dwa niezależne. Smartfony w pobliżu przekazują pozycję znacznika. Gdy ktoś znajdzie gondolę przed zespołem i podejdzie ze smartfonem, zespół zobaczy nową pozycję.
  \item Zapasowo: smartfon z funkcją „Znajdź moje urządzenie” -- w 2020 roku pomógł znaleźć ładunek~\cite{raport2020}.
\end{itemize}
\end{sekcja}

\begin{sekcja}{Zasilanie}
\begin{itemize}
  \item Używaj baterii litowych, np. Energizer Ultimate Lithium (Energizer Extreme). Działają w niskiej temperaturze i są lekkie.
  \item Włóż nowe baterie przed każdym lotem.
  \item Licz budżet energii: pobór każdego urządzenia $\times$ czas pracy $\times$ zapas. Uwzględnij prąd rozruchowy serw i grzałek.
  \item Czas pracy: co najmniej dwa razy dłuższy niż planowany lot plus czas poszukiwań.
\end{itemize}
\end{sekcja}

\begin{sekcja}{Termika: mróz i próżnia}
W stratosferze powietrze ma do \SI{-60}{\degreeCelsius}, ale jest bardzo rzadkie. Prawie nie ma konwekcji, więc elektronika \textbf{nie oddaje ciepła} do powietrza. Słońce dodatkowo nagrzewa gondolę. W misji z lipca 2020 roku elektronika w izolowanej gondoli przekroczyła \SI{50}{\degreeCelsius}~\cite{praca2022}.

\textbf{Zasady projektowe:}
\begin{itemize}
  \item projektuj elektronikę o wysokiej sprawności -- mniej strat to mniej ciepła,
  \item dodaj duży zapas pojemności cieplnej (masa, która akumuluje ciepło, np. baterie, metalowa płyta),
  \item odprowadzaj ciepło przez przewodzenie do obudowy, nie przez powietrze,
  \item nie dodawaj izolacji ani grzałek bez bilansu cieplnego dla startu (zimno) i dla stratosfery (brak chłodzenia).
\end{itemize}

\begin{wskazowka}Do testów w niskiej temperaturze zapytaj biologów na Wydziale Chemicznym o zamrażarkę laboratoryjną (\SI{-80}{\degreeCelsius}).\end{wskazowka}
\end{sekcja}
